% ============================================================
%  preamble.tex — 日本語翻訳版プリアンブル
%  原本の mystyle.sty + notations.sty + 日本語環境を統合
% ============================================================

% --- 基本数学パッケージ ---
\usepackage{amsmath,amsfonts,amssymb,amsthm}
\usepackage{mathrsfs}
\usepackage{mathtools}
\usepackage{stmaryrd}
\usepackage{dsfont}
\usepackage{bbold}

% --- 日本語パッケージ ---
\usepackage[dvipdfmx,hidelinks]{hyperref}
\usepackage{pxjahyper}

% --- 図表 ---
\usepackage[dvipdfmx]{graphicx}
\graphicspath{{./images/}}
\usepackage{subfigure}
\usepackage{array}
\usepackage{tabularx}
\usepackage{xcolor}

% --- リスト ---
\usepackage{enumerate}
\usepackage{enumitem}

% --- float ---
\usepackage{float}
\floatstyle{boxed}
\newfloat{listing}{H}{loi}
\floatname{listing}{Table}

% --- 灰色ボックス (mdframed) ---
\usepackage{mdframed}
\mdfsetup{backgroundcolor=gray!5}
\mdfdefinestyle{mystyle}{backgroundcolor=gray!15}

% ============================================================
%  定理環境（原本と同じ環境名、日本語ラベル）
% ============================================================
\newtheorem{thm}{定理}[chapter]
\newtheorem{prop}{命題}[chapter]
\newtheorem{rem}{注意}[chapter]
\newtheorem{defn}{定義}[chapter]
\newtheorem{cor}{系}[chapter]
\newtheorem{lem}{補題}[chapter]
\newtheorem{exmp}{例}[chapter]

% --- 環境エイリアス（一部の章で長い名前が使われるため）---
\newenvironment{proposition}[1][]{\begin{prop}[#1]}{\end{prop}}
\newenvironment{example}[1][]{\begin{exmp}[#1]}{\end{exmp}}
\newenvironment{remark}[1][]{\begin{rem}[#1]}{\end{rem}}

% --- 灰色ボックス付き注意環境 (rem1: 薄灰, rem2: 濃灰) ---
\newenvironment{rem1}[1]{%
\begin{rem}[#1]\begin{mdframed}
}{%
\end{mdframed}\end{rem}}

\newenvironment{rem11}[2]{%
\begin{rem}[#1]\label{#2}\begin{mdframed}
}{%
\end{mdframed}\end{rem}}

\newenvironment{rem2}[1]{%
\begin{rem}[#1]\begin{mdframed}[style=mystyle]
}{%
\end{mdframed}\end{rem}}

% --- rs リスト環境 ---
\newenvironment{rs}{\begin{list}{\rule[0.6ex]{1.3ex}{0.18ex}~}{
 \topsep=0.3ex \itemsep=0.3ex \labelsep=0em \parsep=0em
 \listparindent=1em \itemindent=0em
 \settowidth{\labelwidth}{--~} \leftmargin=\labelwidth
}}{\end{list}}

% ============================================================
%  文字マクロ（mystyle.sty より）
% ============================================================

%% 二重線文字
\newcommand{\NN}{\mathbb{N}}
\newcommand{\CC}{\mathbb{C}}
\newcommand{\GG}{\mathbb{G}}
\newcommand{\LL}{\mathbb{L}}
\newcommand{\PP}{\mathbb{P}}
\renewcommand{\SS}{\mathbb{S}}
\newcommand{\QQ}{\mathbb{Q}}
\newcommand{\RR}{\mathbb{R}}
\newcommand{\VV}{\mathbb{V}}
\newcommand{\ZZ}{\mathbb{Z}}
\newcommand{\FF}{\mathbb{F}}
\newcommand{\KK}{\mathbb{K}}
\newcommand{\TT}{\mathbb{T}}
\newcommand{\UU}{\mathbb{U}}
\newcommand{\EE}{\mathbb{E}}
\newcommand{\XX}{\mathbb{X}}
\newcommand{\YY}{\mathbb{Y}}

%% カリグラフィ文字
\newcommand{\Aa}{\mathcal{A}}
\newcommand{\Bb}{\mathcal{B}}
\newcommand{\Cc}{\mathcal{C}}
\newcommand{\Dd}{\mathcal{D}}
\newcommand{\Ee}{\mathcal{E}}
\newcommand{\Ff}{\mathcal{F}}
\newcommand{\Gg}{\mathcal{G}}
\newcommand{\Hh}{\mathcal{H}}
\newcommand{\Ii}{\mathcal{I}}
\newcommand{\Jj}{\mathcal{J}}
\newcommand{\Kk}{\mathcal{K}}
\newcommand{\Ll}{\mathcal{L}}
\newcommand{\Mm}{\mathcal{M}}
\newcommand{\Nn}{\mathcal{N}}
\newcommand{\Oo}{\mathcal{O}}
\newcommand{\Pp}{\mathcal{P}}
\newcommand{\Qq}{\mathcal{Q}}
\newcommand{\Rr}{\mathcal{R}}
\newcommand{\Ss}{\mathcal{S}}
\newcommand{\Tt}{\mathcal{T}}
\newcommand{\Uu}{\mathcal{U}}
\newcommand{\Vv}{\mathcal{V}}
\newcommand{\Ww}{\mathcal{W}}
\newcommand{\Xx}{\mathcal{X}}
\newcommand{\Yy}{\mathcal{Y}}
\newcommand{\Zz}{\mathcal{Z}}

%% ギリシャ文字略記
\newcommand{\al}{\alpha}
\newcommand{\la}{\lambda}
\newcommand{\ga}{\gamma}
\newcommand{\Ga}{\Gamma}
\newcommand{\La}{\Lambda}
\newcommand{\si}{\sigma}
\newcommand{\Si}{\Sigma}
\newcommand{\be}{\beta}
\newcommand{\de}{\delta}
\newcommand{\De}{\Delta}
\let\tphi\phi
\renewcommand{\phi}{\varphi}
\renewcommand{\th}{\theta}
\newcommand{\om}{\omega}
\newcommand{\Om}{\Omega}

% ============================================================
%  数学演算子（mystyle.sty より）
% ============================================================

\newcommand{\guill}[1]{``#1''}
\renewcommand{\epsilon}{\varepsilon}
\newcommand{\Id}{\mathrm{Id}}
\newcommand{\eqdef}{\ensuremath{\stackrel{\mbox{\upshape\tiny def.}}{=}}}
\DeclareMathOperator{\Ker}{Ker}
\DeclareMathOperator{\tr}{tr}
\DeclareMathOperator{\trace}{tr}
\DeclareMathOperator{\Tr}{Tr}
\DeclareMathOperator{\Supp}{Supp}
\DeclareMathOperator{\Sign}{Sign}
\DeclareMathOperator{\sign}{sign}
\DeclareMathOperator{\support}{supp}
\DeclareMathOperator{\diverg}{div}
\DeclareMathOperator{\diag}{diag}
\DeclareMathOperator{\dom}{dom}
\newcommand{\foralls}{\forall \,}
\renewcommand{\d}{\mathrm{d}}
\newcommand{\pd}[2]{ \frac{ \partial #1}{\partial #2} }
\newcommand{\pdd}[2]{ \frac{ \partial^2 #1}{\partial #2^2} }
\newcommand{\dotp}[2]{\langle #1,\,#2\rangle}
\newcommand{\norm}[1]{\left\| #1 \right\|}
\newcommand{\norms}[1]{\| #1 \|}
\newcommand{\normb}[1]{\Big|\!\Big| #1 \Big |\!\Big|}
\newcommand{\normB}[1]{\left|\!\left| #1 \right|\!\right|}
\newcommand{\abs}[1]{\left\lvert#1\right\rvert}
\newcommand{\absb}[1]{\Big| #1 \Big|}
\newcommand{\func}[4]{ {\left\{  \begin{array}{ccc} #1 & \longrightarrow & #2 \\ #3 & \longmapsto & #4 \end{array}  \right.} }
\newcommand{\transp}[1]{ {#1}^{\mathrm{T}} }
\newcommand{\eq}[1]{\begin{equation*}#1\end{equation*}}
\newcommand{\eql}[1]{\begin{equation}#1\end{equation}}
\DeclareMathOperator{\argmin}{argmin}
\DeclareMathOperator{\argmax}{argmax}
\newcommand{\uargmin}[1]{\underset{#1}{\argmin}\;}
\newcommand{\uargmax}[1]{\underset{#1}{\argmax}\;}
\newcommand{\umin}[1]{\underset{#1}{\min}\;}
\newcommand{\umax}[1]{\underset{#1}{\max}\;}
\newcommand{\usup}[1]{\underset{#1}{\sup}\;}
\newcommand{\uinf}[1]{\underset{#1}{\inf}\;}
\newcommand{\pa}[1]{\left( #1 \right)}
\newcommand{\bpa}[1]{\big( #1 \big)}
\newcommand{\choice}[1]{ %
	\left\{ %
		\begin{array}{l} #1 \end{array} %
	\right. }
\newcommand{\ens}[1]{ \{ #1 \} }
\newcommand{\enscond}[2]{ \left\{ #1 \;:\; #2 \right\} }
\newcommand{\ensconds}[2]{ \{ #1 \;:\; #2 \} }
\newcommand{\range}[1]{\llbracket#1\rrbracket}

% ============================================================
%  スペーシングマクロ（mystyle.sty より）
% ============================================================
\newcommand{\qandq}{ \quad \text{and} \quad }
\newcommand{\qqandqq}{ \qquad \text{and} \qquad }
\newcommand{\qifq}{ \quad \text{if} \quad }
\newcommand{\qqifqq}{ \qquad \text{if} \qquad }
\newcommand{\qwhereq}{ \quad \text{where} \quad }
\newcommand{\qqwhereqq}{ \qquad \text{where} \qquad }
\newcommand{\qwithq}{ \quad \text{with} \quad }
\newcommand{\qqwithqq}{ \qquad \text{with} \qquad }
\newcommand{\qforq}{ \quad \text{for} \quad }
\newcommand{\qqforqq}{ \qquad \text{for} \qquad }
\newcommand{\qqsinceqq}{ \qquad \text{since} \qquad }
\newcommand{\qsinceq}{ \quad \text{since} \quad }
\newcommand{\qarrq}{\quad\Longrightarrow\quad}
\newcommand{\qqarrqq}{\quad\Longrightarrow\quad}
\newcommand{\qiffq}{\quad\Longleftrightarrow\quad}
\newcommand{\qqiffqq}{\qquad\Longleftrightarrow\qquad}
\newcommand{\qsubjq}{ \quad \text{subject to} \quad }
\newcommand{\qqsubjqq}{ \qquad \text{subject to} \qquad }
\newcommand{\qstq}{ \quad \text{s.\ t.\ }\quad }

% ============================================================
%  図マクロ（mystyle.sty より）
% ============================================================
\newcommand{\myfigure}[3]{
    \begin{figure}[ht]
    \begin{center}
        #1
    \end{center}
        \caption{\textit{#2}}
    \label{#3}
    \end{figure}
}
\newcommand{\myfigurestar}[3]{
    \begin{figure*}[ht]
    \begin{center}
        #1
    \end{center}
    \vspace{-4mm}
        \caption{\textit{#2}}
    \label{#3}
    \end{figure*}
}
\newcommand{\imgBox}[1]{%
\setlength{\fboxsep}{0pt}
\fbox{#1}}
\newcommand{\figitem}[1]{\textsf{\textbf{(#1)}}}
\newcommand{\figitemraise}[2]{\raisebox{#1}{\figitem{#2}}}
\newcommand{\itemraise}[2]{\raisebox{#1}{#2}}

% ============================================================
%  OT固有記法（notations.sty）
% ============================================================
\usepackage{notations}

% ============================================================
%  i.e. / e.g. の再定義
% ============================================================
\def\ie{すなわち}
\def\eg{例えば}

% --- 参考文献 ---
\usepackage[numbers,sort&compress]{natbib}
